% 18_body.tex — Day4 산출물 3/5 (⑱) 종료 보고서 · 미해결 사항 이관 목록 · 계약 종결 (빈 템플릿 / 완성 예시 공용 본문)
% 드라이버: 18_종료보고서_빈템플릿.tex (\wbblanktrue) / 18_종료보고서_완성예시.tex (\wbblankfalse)
% 내용 출처: PM트랙/Day4_워크북.md §3.3(항목 가이드) §3.4(빈 템플릿) §3.5(온담 P1 완성 예시본)
\documentclass[10pt,a4paper]{article}
\usepackage[a4paper,margin=16mm,top=14mm,bottom=20mm]{geometry}
\def\DocNum{3}
\def\DocTitle{종료 보고서 · 미해결 사항 이관 · 계약 종결}
\def\DocSub{⑱ End Project Report, Follow-on Actions \& Contract Closure}
\def\DocVariant{\B{빈 템플릿}{온담 P1 완성 예시}}
\usepackage{wbstyle}
\begin{document}

\docheader
 {\Bg{[정식 명칭과 코드]}{차세대 주문·재고 통합 플랫폼 구축 (ONE ONDAM P1)}}
 {\Bg{[데이터 일자(G4) / 발행 / 보고]}{데이터 2027-05-28 (G4) / 발행 06-04 / 이사회 06월}}
 {\Bg{[P1 PM]}{P1 PM\rd{7mm}}}
 {\Bg{[스폰서 실명 — G4 승인]}{정미란 CFO 05-28 (G4 승인)\rd{7mm}}}

% ------------------------------------------------------------------ 1
\secbar{1. 문서 정보 · 완료 기준 판정}
\guide{보고 대상·시점, 기간과 기준선 이력(\textbf{재기준선 횟수} — 한 번도 없었다는 사실이 첫 문장이다), 판정자·승인자, 입력 문서. 헌장 완료 기준을 \textbf{문언 그대로} 인용하고 옆에 판정(충족/조건부)·근거·판정자를 적는다. G4 시점에 기간이 안 찬 기준(가용성 3개월)은 “G4 시점 실적 + 잔여 기간 판정자”로. “성공적으로 종료”는 판정이 아니다.}
{\footnotesize
\begin{tabularx}{\linewidth}{|L{28mm}|Y|}\hline
\hd{항목} & \hd{내용} \\ \hline
\B{\rd{9mm}}{}기간 / 기준선 & \Bg{[착수 ~ 종료(n개월) / BL-01(일자) → BL-02(일자), 재기준선 n회]}{2026-01-05 ~ 2027-05-28(17개월, 헌장 05-29 토요일 → 직전 영업일) / BL-01(G1 04-30) → BL-02(10-23), \textbf{재기준선 0회}} \\ \hline
\B{\rd{9mm}}{}승인 & \Bg{[판정(실명 — SAC / 형상·문서) / 승인 실명·일자 / 이사회 보고]}{판정 박준영(SAC)·박세연(형상·문서) / 승인 \textbf{정미란} 05-28 / 이사회 보고 김선호 의장 06월 정기} \\ \hline
\B{\rd{10mm}}{}입력 & \Bg{[성과보고 최종호, 변경 로그, 등록부, 품질 기록, 계약 종결 서면, SAC 판정, ELS 종료 보고, 편익 원장]}{성과보고 제17호(05-31), 변경 로그 v1.4, 등록부 v1.4, 품질 통제 기록 최종, 계약 종결 서면 7건, SAC 판정(05-21), ELS 종료 보고(04-23), 편익 실현 원장 v1.0(⑳)} \\ \hline
\B{\rd{8mm}}{}사본 & \Bg{[내부감사 등]}{윤태호 내부감사팀장, 프로그램 PMO장} \\ \hline
\end{tabularx}}
\vspace{1.5mm}
{\footnotesize
\begin{tabularx}{\linewidth}{|Y|L{36mm}|L{58mm}|L{22mm}|}\hline
\hd{완료 기준 (헌장 \B{§n 문언}{v1.1 §14 정상 종료})} & \hd{판정} & \hd{근거 · 일자} & \hd{판정자} \\ \hline
\ifwbblank
\rd{9mm}\g{① [SAC 미해소 0건 또는 해소 기한·책임자 서명된 목록]} & \g{[충족 / 조건부 — n + n]} & & \g{[인수자]} \\ \hline
\rd{9mm}\g{② [소스·형상 인수 100\%]} & & & \\ \hline
\rd{9mm}\g{③ [인터페이스 문서 인수]} & & & \\ \hline
\rd{9mm}\g{④ [컷오버 후 n개월 가용성 n\%]} & \g{[G4 시점 실적 + 잔여 판정자]} & & \\ \hline
\rd{9mm}\g{⑤ [종료 보고서·교훈 제출]} & & & \\ \hline
\else
① SAC 미해소 0건 또는 해소 기한·책임자 서명된 목록 & \textbf{충족} (11 충족 + 1 조건부 서명) & SAC 판정 05-21, 조건부 = 아카이브 접근 절차 재경 승인 06-15(재경팀장) & 박준영 \\ \hline
② 소스·형상 발주사 인수 100\% & 충족 & 클린 빌드 재현 100\%(05-12, 하도급 220 FP 포함), SBOM 312건, 빌드·배포 절차서 & 박세연·형상 담당 \\ \hline
③ 인터페이스 47본 설계문서 인수 & 충족 & 47/47(표준 양식, AF-1-02 9본 보완 10-30) & 박세연 \\ \hline
④ 컷오버 후 3개월 가용성 99.5\% & 충족 (03월 99.872\% / 04월 100\% / 05-01~28 100\%) & APM 가용성 보고, 계획 정지 제외 & 박준영 \\ \hline
⑤ 종료 보고서·교훈 기록 제출 & 충족 & 본 문서·교훈 등록부(⑲) 06-04 & P1 PM \\ \hline
\fi
\end{tabularx}}

% ------------------------------------------------------------------ 2 (가로) 최종 성과표
\landscapeon
\secbar{2. 기준선 대비 최종 성과 \B{}{(BL-02 대비 — day4\_closing.py §2·3·7)}}
\guide{축별로 기준선(BL-0n)·최종 실적·편차·허용범위(격자)·판정, BL-01 대비는 참고 열로만. 성과보고 최종호의 표를 그대로 가져온다. 승인 예산 대비 “n억 절감”으로 적지 않는다 — 관리예비비와 반납 컨틴전시는 절감이 아니라 \textbf{미사용}이다. 최종 CPI를 실행 중 조정 CPI와 나란히 놓아 “그때 본 것이 맞았는가”를 보인다.}
{\footnotesize\setlength{\tabcolsep}{2.5pt}
\begin{tabularx}{\linewidth}{|L{16mm}|L{42mm}|Y|L{22mm}|L{30mm}|L{44mm}|L{24mm}|}\hline
\hd{축} & \hd{기준선 \B{(BL-0n)}{BL-02}} & \hd{최종 실적 \B{(G4)}{(05-28)}} & \hd{편차} & \hd{허용범위 (격자)} & \hd{판정} & \hd{참고 BL-01} \\ \hline
\ifwbblank
\rd{10mm}범위 & \g{[FP / RTM / WBS]} & \g{[인수 FP (n\%), RTM 검수 연결, 산출물 n종]} & & \g{[±FP / ±건]} & & \\ \hline
\rd{10mm}일정 & \g{[종료 게이트 일자, 컷오버 창]} & \g{[실제 종료일, 창 정시 여부]} & \g{[n영업일]} & \g{[+n영업일 / 창 0]} & \g{[이내 — 실행 중 전망 대비]} & \\ \hline
\rd{10mm}원가 & \g{[PMB / 집행 한도 / 승인]} & \g{[AC / CPI]} & \g{[VAC / 한도 여유]} & \g{[EAC > 한도 예외]} & \g{[예외 보고 n회 · 경고 n회]} & \\ \hline
\rd{10mm}품질 & \g{[결함 밀도 기준, 심각도 1·2 미해결 0]} & \g{[n건 → n건/FP; 미해결]} & \g{[기준의 n\%]} & & & \\ \hline
\rd{10mm}리스크 & \g{[잔여 EMV / 단일]} & \g{[종결·이관, 최대 잔여, 실현 n건]} & & & & \\ \hline
\rd{10mm}편익 & \g{[전망 n억, $-$n\%p]} & \g{[G4 시점 첫 측정]} & \g{[미판정]} & & \g{[원장으로 이관]} & \\ \hline
\rd{10mm}지속가능성 & \g{[High 미조치 0, 사고 0]} & & & & & \\ \hline
\else
범위 & FP 12,585 / RTM 1,189 / WBS 43 & 인수 FP \textbf{12,585 (100\%)}, RTM 검수 연결 100\%, 산출물 5종 인수 & 0 & ±372 FP / ±35건 & 이내 & 12,400 → +185(+1.5\%) \\ \hline
일정 & G4 2027-05-28, 컷오버 창 02-26·(2차) 03-27 & \textbf{G4 05-28 정시}, 1차 창 정시(34h), 2차 창 정시(8h) & 0영업일 & +15영업일(06-18) / 창 0 & \textbf{이내} — 9월 as-is 전망 06-11(+10)·ES 추세 07-12(+31)를 EX-01 대안 1로 회복 & 동일 \\ \hline
원가 & PMB \textbf{146.82} / 집행 한도 156.00 / 승인 168.00 & \textbf{AC 151.95} / CPI \textbf{0.966} & VAC \textbf{$-$5.13} / 집행 한도 여유 \textbf{+4.05} & EAC > 156.00 예외 & \textbf{이내} — 예외 보고 0회(경고 2회: 10월·1월) & 145.80 → +6.15 \\ \hline
품질 & 0.4건/FP(통합시험~UAT ÷ 12,585), 심각도 1·2 미해결 0(G3) & 통합시험 612 + UAT 575 = \textbf{1,187건 → 0.094건/FP}; 심각도 1·2 미해결 G3 0·G4 0; 총 발견 2,563(설계 210·스프린트 1,166·통합~UAT 1,187) & 기준의 24\% & 0.4 & 이내 (9월 전망 0.09 적중) & — \\ \hline
리스크 & 잔여 EMV 8억 / 단일 2억 & G4 종결·이관, 최대 잔여(9월) 7.26, 실현 4건(R-04·05·10 부분·R-14 미실현) & — & 8 / 2 & 이내 (여유 최소 0.74, 11월) & — \\ \hline
편익 & G2 전망 118억, $-$5\%p & G4 시점 첫 측정: B-08 13.6h(목표 12h, M+12 판정), 두 분모 리포트 가동 — 판정은 G5(2028-03-31) & 미판정 & $-$5\%p & ⑳ 원장으로 이관 & — \\ \hline
지속가능성 & High 미조치 0, 사고 0 & 보안 High 0(2회차 01-20), 개인정보 사건 1건(P1-02, 신고·복구 완료) & — & 0 & 사고 1건 기록 — 조치 완료, 규제 제재 없음(06-04 현재) & — \\ \hline
\fi
\end{tabularx}}

\needspace{60mm}\vspace{2mm}\noindent\textbf{\small 원가 항목별} \g{(억 원, BL-02 PV 대비 — CV의 원인과 자금원, 실행 중 ETC 대비)}\par\vspace{0.5mm}
{\footnotesize\setlength{\tabcolsep}{2.5pt}
\begin{tabularx}{\linewidth}{|L{22mm}|C{18mm}|C{18mm}|C{16mm}|Y|L{50mm}|}\hline
\hd{항목} & \hd{PV \B{(BL)}{(BL-02)}} & \hd{AC 최종} & \hd{CV} & \hd{원인 (자금원)} & \hd{\B{제n호}{제9호} ETC 대비} \\ \hline
\ifwbblank
\rd{7mm}SI & & & & \g{[컨틴전시 집행 n건 — 고정가는 CV = 컨틴전시 소진]} & \g{[ETC n → 실 n (±)]} \\ \hline
\rd{7mm}상용SW & & & & & \\ \hline
\rd{7mm}인프라 & & & & \g{[실사용 편차]} & \\ \hline
\rd{7mm}단말 & & & & & \\ \hline
\rd{7mm}데이터 전환 & & & & & \\ \hline
\rd{7mm}시험·감리 & & & & & \\ \hline
\rd{7mm}PMO & & & & & \\ \hline
\rd{7mm}\textbf{합계} & & & & & \g{[EAC4 n → ±n (±2 내?)]} \\ \hline
\else
SI & 69.22 & 71.02 & $-$1.80 & 컨틴전시 집행 6건: R-07 0.15·R-05 0.42·R-03 0.50·R-19 인력 0.35·R-20 대기 0.18·R-13 패치 0.20 (고정가 — CV = 컨틴전시 소진) & ETC 39.71 → 실 40.14(+0.43: R-13·R-20·R-19 인력분) \\ \hline
상용SW & 19.40 & 19.42 & $-$0.02 & 라이선스 잔금 4.70 01-15 발생(통합시험 통과), 임시 라이선스 0.02(미배정 풀) & 5.90 → 5.92 \\ \hline
인프라 & 20.90 & 22.16 & $-$1.26 & 클라우드 실사용 +0.76(6×0.5 + 8×0.42 $-$ 14×0.4), 망분리·보안 +0.40, DR +0.10 — 유보 3.0 미사용 & 7.92 → 7.66 \\ \hline
단말 & 8.60 & 8.90 & $-$0.30 & 단일 공급사 견적(09-18), 미배정 풀 & 8.90 → 8.90 \\ \hline
데이터 전환 & 9.40 & 10.05 & $-$0.65 & 아카이브 정책 0.20(미배정 풀), R-10 재정제 0.45(컨틴전시) & 5.10 → 5.75(+0.65) \\ \hline
시험·감리 & 9.70 & 10.20 & $-$0.50 & 보안 1회차 +0.05, 부하 시나리오 +0.20, R-19 회귀 도구 0.25 & 7.40 → 7.65 \\ \hline
PMO & 9.60 & 10.20 & $-$0.60 & 월 0.60 × 17(도구·회의) & 4.80 → 4.80 \\ \hline
\textbf{합계} & \textbf{146.82} & \textbf{151.95} & \textbf{$-$5.13} & & EAC4 150.86 → \textbf{+1.09}(±2 내), EAC1′ 150.57, EAC3′ 155.04 미도달 \\ \hline
\fi
\end{tabularx}}
\B{\small\g{EAC 이력(호별 공식 EAC, 경고 시점) / 실행 중 조정 CPI vs 최종 CPI / 예외 보고 횟수: }\hrulefill\par}{\small \textbf{EAC 이력}(공식 EAC4): 제9호(09-30) 150.86 → 제10호 151.20 → 제11호(11-30) 152.40(R-10 재정제·CR-13, \textbf{경고 A}) → 제13호(01-31) 151.70 → 제15호(03-31) 152.10(임시 라이선스·ELS 인력, 경고) → 제17호(05-31) \textbf{151.95 확정}. 9월 조정 CPI 0.968 → 최종 CPI 0.966 — 발생주의 조정값이 최종값을 0.002 안에서 예측했다. 예외 보고 EX-01(일정) 1회, 원가 예외 0회.\par}

% ------------------------------------------------------------------ 3 (가로) 변경 로그 총괄
\clearpage
\secbar{3. 변경 로그 총괄 \B{CR-01 ~}{CR-01 ~ CR-14 (CR-01~11은 Day3 §1.5 항목 10 요약)}}
\guide{전 변경의 한 줄 요약 — 접수·제목·FP Δ·금액·판정 경로·상태·BL. \textbf{기각·이관된 변경도 남긴다} — 후속 과제의 출발점이다. 집계 행에 경로별 건수, FP 누적이 허용범위 ±의 몇 \%인지, 재확인 절차 트리거 발동 여부를 적는다.}
{\scriptsize\setlength{\tabcolsep}{2pt}
\begin{longtable}{|C{12mm}|L{13mm}|L{76mm}|C{13mm}|C{22mm}|L{52mm}|L{46mm}|L{16mm}|}\hline
\hd{CR} & \hd{접수} & \hd{제목} & \hd{FP Δ} & \hd{금액 (억)} & \hd{판정 · 경로} & \hd{상태} & \hd{BL} \\ \hline \endhead
\ifwbblank
\rd{6mm}CR-01 & & & & & \g{[PM 전결 / CCB / 조건부 / 기각 / 이관]} & \g{[완료 — 후속 사건]} & \g{[BL-0n]} \\ \hline
\rd{6mm}CR-02 & & & & & & & \\ \hline
\rd{6mm}CR-03 & & & & & & & \\ \hline
\rd{6mm}CR-04 & & & & & & & \\ \hline
\rd{6mm}CR-05 & & & & & & & \\ \hline
\rd{6mm}CR-06 & & & & & & & \\ \hline
\rd{6mm}CR-07 & & & & & & & \\ \hline
\rd{6mm}CR-08 & & & & & & & \\ \hline
\rd{6mm}CR-09 & & & & & & & \\ \hline
\rd{6mm}CR-10 & & & & & & & \\ \hline
\rd{6mm}CR-11 & & & & & & & \\ \hline
\rd{6mm}CR-12 & & & & & & & \\ \hline
\rd{6mm}CR-13 & & & & & & & \\ \hline
\rd{6mm}CR-14 & & & & & & & \\ \hline
\rd{6mm} & & & & & & & \\ \hline
\rd{8mm}\textbf{집계} & & \g{[n건: 승인 n · 조건부 n · 거버넌스 n · 기각 n · 이관 n]} & \g{[+n (n\%)]} & \g{[+n]} & \g{[허용 ±n의 n\%, 재확인 트리거(±n\%) 발동 여부]} & & \\ \hline
\else
CR-01 & 02-11 & 배달3사 IF 규격 v2 & 0 & 0 & PM 전결 & 완료 & BL-01 \\ \hline
CR-02 & 02-26 & 재고 정확도 두 분모 리포트(StRS-034) & +15 & 0.08 & PM 전결(기준선 내) & 완료 — \textbf{12-18 첫 리포트} & BL-01 \\ \hline
CR-03 & 03-11 & 엑셀 발주 양식 5종 & +22 & 0.12 & PM 전결·CCB 사후 & 완료 & BL-02 \\ \hline
CR-04 & 03-25 & 옴니 고도화 3종 R5 이연 & 0 & 0 & CCB & 완료 — \textbf{R5 05-07 릴리스} & BL-01 \\ \hline
CR-05 & 03-25 & 배달 주문 가맹 배분·정산(StRS-054) & 0 & 0 & Won't & 기각(기록) & — \\ \hline
CR-06 & 06-10 & FZ-05 승인 16건 & +58 & 0.32 & 변경협의체 → PM·CCB 사후 & 완료 & BL-02 \\ \hline
CR-07 & 09-17 & 인터페이스 목록 갱신(+2 $-$2) & +40 & 0.22 & PM 전결·CCB 사후 & 완료 & BL-02 \\ \hline
CR-08 & 09-16 & 파일럿 임시 연동 & +90 & 0.495 & CCB(조건 4) & 완료 — 폐기 04-27 & BL-02 \\ \hline
CR-09 & 09-17 & 발주 마감 22:00 파라미터 & +35 & 0.19 & CCB 조건부 & 완료 — 운용 20:00 확정(01-21) & BL-02 \\ \hline
CR-10 & 09-18 & 원가 허용범위 정의 & 0 & 0 & CCB + 스폰서 & 완료(헌장 v1.1 확정) & 헌장 \\ \hline
CR-11 & 09-22 & 전환 5년 한정·아카이브 & $-$60 & $-$0.33(+용역 $-$2.0) & CCB & 완료 — 아카이브 절차 승인 대기(이관 1) & BL-02 \\ \hline
\textbf{CR-12} & 10-20 & 배달 주문 가맹 배분·정산 재상정(오정근, 개산 +160) & 0 & 0 & CCB 11-19: \textbf{P1 범위 외 → 후속 과제 이관}(P3 부속서 협상 결과 후) & 이관 & — \\ \hline
\textbf{CR-13} & 12-16 & 델리 서브 브랜드 상품 코드 체계(R-10 부분 실현, 정본 판정 후 접수) & 0 & 0(재정제 0.45 컨틴전시 R-10) & PM 전결·CCB 사후(12-17) — 데이터 정제 범위, FP 불변 & 완료 & BL-02 \\ \hline
\textbf{CR-14} & 01-12 & 미동의 35점 컷오버 후 2차 배포·구형 단말 펌웨어 패치(IS-06·R-13) & 0 & 0(패치 0.20 컨틴전시 R-13) & CCB 01-21 & 완료 — 잔여 9점(이관 2) & BL-02 \\ \hline
\textbf{집계} & & 14건: 승인 10 · 조건부 1(CR-09) · 거버넌스 1(CR-10) · 기각 1(CR-05) · 이관 1(CR-12) & \textbf{+185 (+1.5\%)} & \textbf{+1.02} SI / 용역 $-$2.0 & 허용 ±372의 49.7\%, 재확인 트리거(±3\%) 미발동 & & BL-02 최종 \\ \hline
\fi
\end{longtable}}

% ------------------------------------------------------------------ 4 컨틴전시 · 관리예비비 총괄
\secbar[80mm]{4. 컨틴전시 · 관리예비비 집행 총괄 \B{}{(day4\_closing.py §3)}}
\guide{컨틴전시 총액이 한 줄로 검산되게 적는다 — \textbf{편입 + 리스크 배정 집행 + 미배정 풀 집행 + 편차 흡수 + 반납 = 총액}. “편차 흡수”(리스크 배정 없이 AC로 나타난 초과)를 따로 적어 컨틴전시가 어디로 갔는지 보인다. 집행별 결재·CCB 보고 이력(누락이 있었으면 소급 사실까지). 관리예비비의 요청·집행·반납, 유보 항목의 정산. “잔여 n억 절감”이라 적지 않는다.}
{\footnotesize\setlength{\tabcolsep}{2.5pt}
\begin{tabularx}{\linewidth}{|L{40mm}|C{22mm}|Y|L{60mm}|}\hline
\hd{구분} & \hd{금액} & \hd{내역} & \hd{결재 · 보고} \\ \hline
\ifwbblank
\rd{8mm}기준선 편입 & & \g{[CR-nn 변경 대가 → PMB]} & \g{[CCB 일자, BL 발행]} \\ \hline
\rd{10mm}리스크 배정 집행 & & \g{[R-nn 금액(월) · … — 소급 보고 건 표시]} & \g{[집행 결재 + CCB 보고 규칙]} \\ \hline
\rd{8mm}미배정 풀 집행 & & \g{[건별 금액]} & \\ \hline
\rd{8mm}편차 흡수 (리스크 배정 없는 AC 초과) & & \g{[항목별]} & \g{[성과보고 CV 열]} \\ \hline
\rd{7mm}\textbf{반납} & & \g{[집행 한도 $-$ AC]} & \\ \hline
\rd{7mm}\textbf{합계} & & \g{[검산식]} & \g{[소진율(편입 제외)]} \\ \hline
\rd{8mm}관리예비비 & & \g{[요청 n · 집행 n · 반납 n — 대상 리스크의 실현 여부]} & \g{[스폰서]} \\ \hline
\rd{10mm}유보 항목 정산 & & \g{[유보 n건 — 반납 / 승인 전 / 감액 확정]} & \\ \hline
\else
BL-02 편입 & 1.02 & CR-03·06·07·08·09·11 변경 대가 → PMB & CCB 10-15, BL-02 10-23 \\ \hline
리스크 배정 집행 & 2.50 & R-07 0.15(06) · R-05 0.42(07) · R-03 0.50(08, \textbf{CCB 소급 보고 10-15}) · R-10 0.45(12) · R-19 0.60(11~01) · R-20 0.18(10-14 판정) · R-13 0.20(01) & 집행 결재 P1 PM + 재경 확인, CCB 월 보고(자동 안건 등재 11월부터) \\ \hline
미배정 풀 집행 & 0.77 & 단말 견적 0.30 · 전환 아카이브 0.20 · 부하 시나리오 0.20 · 보안 1회차 0.05 · 임시 라이선스 0.02 & 동일 \\ \hline
편차 흡수 (리스크 배정 없는 AC 초과) & 1.86 & 클라우드 실사용 0.76 · 망분리·DR 0.50 · PMO 단가 0.60 & 성과보고 CV 열에 매월 기록 \\ \hline
\textbf{반납} & \textbf{4.05} & 집행 한도 156.00 $-$ AC 151.95 & 프로그램 PMO장 06월 \\ \hline
\textbf{합계} & \textbf{10.20} & 1.02 + 2.50 + 0.77 + 1.86 + 4.05 & 소진율(편입 제외) 50.3\% \\ \hline
관리예비비 12.0 & 요청 0 · 집행 0 · \textbf{반납 12.0} & 대상 리스크 R-01·R-14·R-17 미실현(R-01은 분리 의결로 흡수) & 스폰서 \\ \hline
유보 3건 (Day2 조정표) & — & 클라우드 3.0 \textbf{반납}(IS-10, 부하시험 후 사이징 확정) / PMO 3.0 → 하자보수기 PMO \textbf{승인 전}(이관 10) / 전환 2.0 \textbf{감액 확정}(CR-11, DT-01) & \\ \hline
\fi
\end{tabularx}}

% ------------------------------------------------------------------ 5 잔여 리스크 이관
\secbar[70mm]{5. 잔여 리스크의 운영 이관 \B{}{(등록부 v1.4, G4 종결)}}
\guide{등록부 최종 상태(종결·실현 후 종결·이관), 운영 리스크로 이관되는 것의 \textbf{새 소유자·검토 주기}, 편익 리스크로 이관되는 것(⑳). 프로젝트 리스크는 G4에 전부 종결하거나 이관한다 — “잔여”로 남기지 않는다. 신규 운영 리스크(프로젝트가 만든 새 의존성)도 여기서 등록한다.}
{\footnotesize\setlength{\tabcolsep}{2.5pt}
\begin{tabularx}{\linewidth}{|L{34mm}|Y|L{44mm}|L{34mm}|L{28mm}|}\hline
\hd{R} & \hd{최종 상태} & \hd{이관처} & \hd{새 소유자} & \hd{검토} \\ \hline
\ifwbblank
\rd{8mm}\g{[R-nn (묶음 가능)]} & \g{[종결 — 근거]} & \g{—} & \g{—} & \g{—} \\ \hline
\rd{8mm} & \g{[실현(CR-nn) → 잔여]} & \g{[운영 리스크 → KE-nn]} & \g{[실명]} & \g{[CAB / 서비스 리뷰]} \\ \hline
\rd{8mm} & & \g{[운영 계약 리스크]} & & \\ \hline
\rd{8mm} & & \g{[편익 리스크(⑳ 항목 6)]} & & \g{[M+12 리뷰]} \\ \hline
\rd{8mm} & & & & \\ \hline
\rd{8mm} & & & & \\ \hline
\rd{10mm}\g{신규 운영 리스크} & \g{[① ② ③ — 프로젝트가 만든 새 의존성]} & \g{[운영 리스크 등록부]} & & \\ \hline
\else
R-01·R-14 & 종결 — 2차 컷오버 03-27 정시, 트리거 미발동 & — & — & — \\ \hline
R-02 & 종결 — 동의 91\%, 잔여 9점은 이관 2번 & 운영(영업본부) & 오정근 & 07-31 \\ \hline
R-05·R-07·R-11·R-15·R-17·R-19·R-20 & 종결(실현 후 종결 R-05, 미실현 종결 나머지) & — & — & — \\ \hline
R-10 & 부분 실현(CR-13) 후 종결 & — & — & — \\ \hline
\textbf{R-13} & 실현(CR-14) → 잔여 9점 & 운영 리스크 → \textbf{KE-01} & 오정근·박준영 & CAB \\ \hline
\textbf{R-16} & 실현(P2-07) → 우회 적용 & 운영 리스크 → \textbf{KE-02} & 한동석·박준영 & 3PL 재협상 06-30 \\ \hline
\textbf{R-18} & 협상 타결(월 3,600만 원) & 운영 계약 리스크(앱 벤더) & 오정근·박세연 & 서비스 리뷰 \\ \hline
\textbf{R-12} & 이관 & \textbf{편익 리스크(⑳ 항목 6)} & 재무본부·한동석·나영선 & M+12 리뷰 \\ \hline
신규 운영 리스크 & ① 3PL 재협상 결렬 시 폴백 영구화(SLA ⑥ 강등 지속) ② SAP ECC 2027-12 EoM(KE-04 확산) ③ IPO 예비심사 중 변경 동결 창 & 운영 리스크 등록부 & 한동석 / 박세연 / 정미란 & 서비스 리뷰 월 1회 \\ \hline
\fi
\end{tabularx}}

% ------------------------------------------------------------------ 6 (가로) 미해결 사항 이관 목록
\secbar[100mm]{6. 미해결 사항 이관 목록 \B{(Follow-on Action Recommendations)}{(Follow-on Action Recommendations — 10건, 승인 6 · 승인 전 4)}}
\guide{8~10건. 내용·\textbf{출처}(어느 문서의 어느 조건에서 왔는가 — SAC 조건부·CCB 조건·SLA 강등·편익 전제)·소유자 \textbf{실명}·기한·승인 여부(승인 = 예산·결정 있음 / 승인 전 = 무엇이·누가·어느 회의체에서 결정해야 하는가)·수신 조직. 소유자가 부서명이거나 기한이 “추후”이면 종료와 함께 사라진다. \textbf{P1 범위 밖이라는 이유로 편익의 전제 조건을 빼지 않는다.} 이사회 보고 시 “승인 전” 항목을 별도로 읽는다.}
{\footnotesize\setlength{\tabcolsep}{2.5pt}
\begin{tabularx}{\linewidth}{|C{6mm}|Y|L{38mm}|L{26mm}|L{22mm}|L{50mm}|L{30mm}|}\hline
\hd{\#} & \hd{내용} & \hd{출처} & \hd{소유자 (실명)} & \hd{기한} & \hd{승인 여부 (승인 / 승인 전 — 무엇이 · 누가)} & \hd{수신 조직} \\ \hline
\ifwbblank
\rd{10mm}1 & \g{[SAC 조건부 항목의 잔여 작업]} & \g{[SAC 항목 n 조건부, CR-nn 조건]} & \g{[실명(운영 실명)]} & \g{[YYYY-MM-DD]} & \g{[승인 — 절차 확정, 승인만 대기]} & \\ \hline
\rd{10mm}2 & & \g{[CR-nn, IS-nn, R-nn]} & & & & \\ \hline
\rd{10mm}3 & & \g{[CCB 조건]} & & & \g{[승인 전 — 무엇이 필요한가]} & \\ \hline
\rd{10mm}4 & & \g{[SLA 매핑 강등 행]} & & & & \\ \hline
\rd{10mm}5 & & & & & & \\ \hline
\rd{10mm}6 & \g{[편익의 전제 조건 — 범위 밖이어도 적는다]} & \g{[편익 원장 분모, R-nn, 헌장 제외 범위]} & & \g{[이사회 부의]} & \g{[승인 전 — 이사회 재검토]} & \\ \hline
\rd{10mm}7 & & & & & & \\ \hline
\rd{10mm}8 & & & & & & \\ \hline
\rd{10mm}9 & & & & & & \\ \hline
\rd{10mm}10 & \g{[유보 항목 복귀 판단]} & & & & & \\ \hline
\else
1 & 아카이브 접근 절차(5년 초과 이력) 재경팀 승인 — 세무 보존 요건 정합 확인 & SAC 항목 11 조건부, CR-11 조건 ④ & \textbf{재경팀장}(운영 박준영) & 2027-06-15 & 승인(절차 확정, 승인만 대기) & 재무본부·P5 \\ \hline
2 & 가맹 미동의 잔여 9점 — 구형 단말 교체 협의 후 2차 배포 완료(KE-01 해소) & CR-14, IS-06, R-13 & \textbf{오정근} & 2027-07-31 & 승인(P3 예산) & 영업본부·P3 \\ \hline
3 & 가맹 발주 마감 22:00 운용 전환 — 물류 야간조 편성(P2 안정화 후)과 함께 CAB 재상정 & CR-09 조건 ①, CCB 01-21 & \textbf{한동석} & 2027-09 CAB & \textbf{승인 전} — 야간조 편성안·3PL 배송 시각 합의 필요(물류본부·운영위) & 물류본부·P2·CAB \\ \hline
4 & 3PL 대성로지스 준실시간 연동 재협상(2027-06 조항) — 타결 시 SLA ⑥ 복귀·KE-02·06 종결, 결렬 시 폴백 영구화 리스크 & IS-05, 계약 변경 로그 7, SLA 매핑 ⑥ & \textbf{한동석} & 2027-06-30 & 승인(물류본부 과제) & 물류본부·P5 \\ \hline
5 & 앱 벤더 유지보수 계약 갱신(월 3,600만 원, API 장애 복구 3시간 UC) 반영 & R-18, UC ⑦ & \textbf{오정근·박세연} & 2027-06-30 & 승인 & 영업본부·IT본부 \\ \hline
6 & 이천센터·안성 원료공장의 입고 검수·검사기록은 P1 범위 밖(종이·엑셀 유지). 재고 정확도 98.5\% 목표는 검사기록 디지털화가 전제 — 2025-11 이사회 반려(34억 SI 견적) 재검토 필요 & 편익 원장 B-02 분모(전체 SKU 84.6\%), R-12, 헌장 §5 제외 범위 & \textbf{나영선 품질안전본부장} & 2027-09 이사회 부의안 & \textbf{승인 전} — 이사회 재검토 필요 & 품질안전본부·프로그램 이사회 \\ \hline
7 & SAP ECC 2027-12 유지보수 종료 대응 — SAP 애드온·IF-SAP 8본 마이그레이션 계획(KE-04 근본 해소 포함) & 회사 정본 §3.2 제외 범위, KE-04 & \textbf{박세연} & 2027-08-31 계획 수립 & \textbf{승인 전} — 별도 과제 예산 & IT본부·프로그램 이사회 \\ \hline
8 & 클라우드 운영 사이징 확정치(월 0.42)·DR 운영비 P5 연 운영비 27.0 반영, 유보 3.0 반납 확인 & IS-10, 컨틴전시 총괄 & \textbf{박준영·재경팀장} & 2027-06-30 & 승인 & P5·재무본부 \\ \hline
9 & ISMS-P 심사(2027-06) 증적 — P1 플랫폼 보안 운영 절차(문서 8종)·파기 배치 로그·P1-02 사건 기록 제출 & 회사 정본 §3.4, SAC 항목 4 & \textbf{최영주} & 2027-06-30 & 승인(P4 소관) & P4·준법지원팀 \\ \hline
10 & 하자보수기 PMO 잔류 1명(2027-06~12) 예산 3.0 중 1.05 집행 — Day2 조정표 “PMO 유보 3.0”의 복귀 판단 & Day2 §4.5 조정표, 컨틴전시 총괄 유보 & \textbf{프로그램 PMO장} & 2027-06-01 운영위 & \textbf{승인 전} — 스폰서 결재(관리예비비 외 프로그램 유보) & 프로그램 PMO·재무본부 \\ \hline
\fi
\end{tabularx}}

% ------------------------------------------------------------------ 7 (가로) 계약 종결
\secbar[60mm]{7. 계약 종결 — 검수 · 정산 · 유보금 · 하자보수 · 지체상금}
\guide{계약별 검수 확인 일자·정산액(산식)·유보금(율·처리·대체 수단)·하자보수(기간·\textbf{범위 = 결함 vs 개선의 경계})·지체상금(금액과 \textbf{근거} — 0은 “없음”이 아니라 순 지연 0의 산식이다)·서면 번호. 하도급 대금 직접 지급 여부. “정산 완료” 한 줄이면 감사에서 “왜 물리지 않았는가”가 지적된다.}
{\scriptsize\setlength{\tabcolsep}{2pt}
\begin{tabularx}{\linewidth}{|L{22mm}|L{30mm}|L{52mm}|L{30mm}|L{44mm}|Y|L{26mm}|}\hline
\hd{계약} & \hd{검수 확인} & \hd{정산액 (억)} & \hd{유보금} & \hd{하자보수} & \hd{지체상금 (금액 · 근거)} & \hd{서면} \\ \hline
\ifwbblank
\rd{16mm}\g{[SI]} & \g{[릴리스별 검수, 최종 검수 확인서 일자]} & \g{[계약 + 변경 대가 + 대기 비용 + 컨틴전시 귀속분 = n — 잔금 지급일]} & \g{[n\% = n억 → 처리(보증증권 대체 등)]} & \g{[기간 / 범위 = 결함, 개선은 유지보수 계약 / 대기 인력·응답 기준]} & \g{[n — 근거 ① 완료일 = 기준선 ② 순 지연 분해 ③ 상계 ④ 동시 지연 원칙]} & \g{[합의서 n호 · 검수 확인서 · 증권]} \\ \hline
\rd{12mm}\g{[데이터 전환 용역]} & & & & & & \\ \hline
\rd{10mm}\g{[상용SW]} & & & & \g{[유지보수 → 운영 계약]} & & \\ \hline
\rd{10mm}\g{[클라우드]} & & & & \g{[운영 계약 승계]} & & \\ \hline
\rd{10mm}\g{[감리·진단·시험]} & & & & & & \\ \hline
\rd{10mm}\g{[단말]} & & & & & & \\ \hline
\rd{10mm}\g{[하도급]} & \g{[수행사 검수 경유, 소스 커밋 확인]} & \g{[직접 지급 요청 유무]} & & \g{[수행사 연대]} & & \\ \hline
\else
SI 한빛디지털 & 릴리스별 R1~R5 검수 완료, 최종 검수 확인서 \textbf{05-26} & 계약 69.22(68.20 + 변경 대가 1.02) + 대기 비용 0.60(R-05 0.42·R-20 0.18) + 컨틴전시 집행 수행사 귀속분 1.20(R-07·R-03·R-19 인력·R-13) = \textbf{71.02} — 잔금 05-31 지급 & \textbf{5\% = 3.46억}(계약 §16 [추가 설정]) → 수행사 선택으로 \textbf{하자보수보증증권 대체}(05-26) & \textbf{1년 2027-05-29 ~ 2028-05-28}, 범위 = 인수 산출물의 결함(심각도 무관), 개선·요구 변경은 유지보수 계약(P5). 대기 인력 4명 → 상주 아님, 응답 UC ② 기준 & \textbf{0} — 근거: ① G4 05-28 = 기준선 완료일 → §18 “지연” 불성립 ② 순 지연 분해(Day3 ⑮ 항목 6 누적): 발주 7 + 동시 3 + 수행 3 $-$ 회복 3 = 10 → EX-01 대안 1로 0 ③ 수행사 단독 3일은 회복 3일과 상계(변경협의체 10-14 확인) ④ 동시 3일은 SCL 원칙 10(공기 연장 인정·비용 불인정) & SI 합의서 6호(10-23)·7호(정산, 05-26), 검수 확인서 R1~R5, 보증증권 \\ \hline
데이터 전환 용역 & 본 전환 검증 6종 통과(02-27), 확인서 03-02 & 9.40 + 아카이브 정책 0.20 + 재정제 0.45 = 10.05 & 3\% [추가 설정] → 보증증권 & 6개월(전환 데이터 정합) & 0 & DT-01(10-23)·DT-02 \\ \hline
상용SW 3사 & 통합시험 통과 01-15(잔금 4.70), API GW 정식 라이선스 03-05 & 19.42(임시 0.02 포함) & — & 라이선스 유지보수 → P5(SW-04) & — & SW-01~05 \\ \hline
클라우드 CSP & 월 청구 & 2년 선약정 6.0 + 종량 누적 6.36 + 기타 = 인프라 22.16 중 & — & P5 운영 계약 승계 & — & CL-01·02 \\ \hline
감리·보안진단·부하시험 & 회차 보고서 접수(G4 감리 05-14) & 10.20 & — & — & — & AU-01·02 \\ \hline
POS 단말 & 입고 검수 12-14, 설치 확인 02-05 & 8.90 & 하자 1년(제조사) & — & 0 & 발주서·검수 조서 \\ \hline
하도급(모빌커넥트, 220 FP) & 수행사 검수 경유, 소스 발주사 저장소 직접 커밋 확인 & 수행사 정산 내 1.21 — \textbf{직접 지급 요청 없음} & — & 수행사 연대 & — & 하도급 승인서 1호 \\ \hline
\fi
\end{tabularx}}
\landscapeoff

% ------------------------------------------------------------------ 8 문서·형상 인수 / 팀 해산 / 승인
\secbar[100mm]{8. 문서 · 형상 인수 / 팀 해산 · 전환 / 승인}
\guide{형상 인수(클린 빌드 재현율·SBOM·라이선스·절차서·정본 전환 시점), 문서 아카이브(형상 저장소 경로·보존 기간), \textbf{감사 지적 번호를 인용}해 “이 절이 그 지적의 답”임을 명시, 팀 해산·전환(사람별 일자 — 해산된 뒤 문서를 찾지 않도록), 하자보수기 체제(창구·판정·보고), 종료 승인 서명. PMO가 하자보수기에 남는데 예산이 없으면 이관 목록에 적는다.}
{\footnotesize
\begin{tabularx}{\linewidth}{|L{28mm}|Y|}\hline
\hd{항목} & \hd{내용} \\ \hline
\B{\rd{14mm}}{}형상 인수 & \Bg{[클린 빌드 재현 n\%(일자, 모듈·하도급 포함), SBOM n건·라이선스, 절차서, 정본 전환 / 감사 지적 번호 대응]}{발주사 저장소 클린 빌드 재현 \textbf{100\%}(05-12, 12개 모듈·하도급 어댑터 포함), SBOM 312 컴포넌트·라이선스 목록(GPL 계열 0), 빌드·배포 절차서 v1.0, 일 1회 동기화 → G4 후 발주사 저장소 단독 정본. \textbf{2023 감사 지적(“위탁 산출물 형상관리·소스 인수 절차 미비”) 대응 완료} — CM-01 계획·AF-1-07 조치·SAC 항목 9} \\ \hline
\B{\rd{12mm}}{}문서 아카이브 & \Bg{[저장소 경로 — 포함 문서 목록, 보존 기간]}{\textsf{governance/p1-closure/} — 헌장 v1.1, 기준선 BL-01·02, RTM v1.4, 성과보고 1~17호, 변경 로그, 등록부 v1.4, 품질 기록, 계약 서면 32건, 런북·SAC, 종료 보고서·교훈·편익 원장. 보존 10년 [추가 설정]} \\ \hline
\B{\rd{12mm}}{}감사 지적 대응 요약 & \Bg{[연도 “지적 문구” → 답이 되는 문서·절]}{2020 “편익 산정 근거 부재” → BC §4 분모·책임자·⑳ 원장 / 2020 “손실 확정 절차 부재” → 헌장 §14 중단 조건·손실 확정(미발동) / 2023 “소스 인수 절차” → 위 형상 인수. 컨틴전시 집행 승인 근거(윤태호 질의) → 항목 4 결재·보고 이력} \\ \hline
\B{\rd{14mm}}{}팀 해산 · 전환 & \Bg{[발주사 n명: 복귀·전환 일자 / 수행사 철수·하자보수 대기 / PMO / PM의 다음 역할]}{발주사 14명: 12명 06-01 원 소속 복귀, 2명 P5 서비스운영팀 전환(06-01) / 수행사: 05-31 철수, 하자보수 대기 4명(비상주) / P1 PMO 3명: 2명 06-01 프로그램 PMO(P2 지원), 1명 하자보수기 잔류(이관 10 승인 전) / P1 PM: 06-01 프로그램 PMO — P2·P5 지원, 하자보수기 창구} \\ \hline
\B{\rd{10mm}}{}하자보수기 체제 & \Bg{[기간 / 창구 → UC / 결함·개선 판정 회의체 / 보고 주기]}{2027-05-29 ~ 2028-05-28. 창구 P5 서비스운영팀(박준영) → 수행사 UC ②. 결함/개선 판정 CAB. 월 1회 서비스 리뷰에 하자 현황 보고} \\ \hline
\B{\rd{10mm}}{}승인 & \Bg{[판정 실명·일자 / 승인 실명·일자 / 이사회 보고 / 사본]}{판정 박준영·박세연 05-21·05-12 / \textbf{승인 정미란 2027-05-28} / 프로그램 이사회 보고 06월(김선호) / 사본 윤태호} \\ \hline
\end{tabularx}}
\vspace{3mm}
\noindent{\footnotesize
\begin{tabularx}{\linewidth}{|Y|Y|Y|Y|}\hline
\hd{작성 — P1 PM (서명)} & \hd{판정 — 운영 인수자 (서명)} & \hd{판정 — 기술 인수자 (서명)} & \hd{승인 — 스폰서 (서명 · 일자)} \\ \hline
\Bg{[P1 PM]}{P1 PM 06-04\rd{10mm}} & \Bg{[서비스운영팀장]}{박준영 05-21} & \Bg{[IT본부장]}{박세연 05-12} & \Bg{[스폰서]}{정미란 2027-05-28 (G4)} \\ \hline
\end{tabularx}}
\end{document}
